\SGANextSectionTarget{sga3:xxiii:subsection:3-4}\SGASubsection{3.4}{Groups of Type \(G_2\)}

\SGARefTarget{sga3:XXIII:3.4.1}
\SGAInlineRefTarget{sga3:xxiii:proposition:3-4-1}\SGAStatement{Proposition 3.4.1}{Let \(S\) be a scheme and \(G\) a
pinned \(S\)-group of type \(G_2\). Write
\[
  \Delta=\{\alpha,\beta\},\qquad
  R^+=\{\alpha,\beta,\alpha+\beta,2\alpha+\beta,
                 3\alpha+\beta,3\alpha+2\beta\}.
\]
\begin{enumerate}
  \item[\SGAInlineRefTarget{sga3:xxiii:proposition:3-4-1:item:i}\SGAInlineRefTarget{sga3:XXIII:3.4.1:item:i}(i)]
    \[
      t_{\alpha\beta}=(w_\alpha w_\beta)^6
        =e=(w_\beta w_\alpha)^6=t_{\beta\alpha}.
    \]
  \item[\SGAInlineRefTarget{sga3:XXIII:3.4.1:item:ii}\SGAInlineRefTarget{sga3:xxiii:proposition:3-4-1:item:ii}(ii)] If
    \[
    \begin{aligned}
      \operatorname{Ad}(w_\beta)X_\alpha&=X_{\alpha+\beta},&
      \operatorname{Ad}(w_\alpha)X_{\alpha+\beta}&=X_{2\alpha+\beta},\\
      \operatorname{Ad}(w_\alpha)X_\beta&=-X_{3\alpha+\beta},&
      \operatorname{Ad}(w_\beta)X_{3\alpha+\beta}&=X_{3\alpha+2\beta},
    \end{aligned}
    \]
    then
    \[
    \begin{aligned}
      \operatorname{Ad}(w_\alpha)X_{2\alpha+\beta}
        &=-X_{\alpha+\beta},&
      \operatorname{Ad}(w_\alpha)X_{3\alpha+\beta}&=X_\beta,\\
      \operatorname{Ad}(w_\alpha)X_{3\alpha+2\beta}
        &=X_{3\alpha+2\beta},&
      \operatorname{Ad}(w_\beta)X_{\alpha+\beta}&=-X_\alpha,\\
      \operatorname{Ad}(w_\beta)X_{2\alpha+\beta}
        &=X_{2\alpha+\beta},&
      \operatorname{Ad}(w_\beta)X_{3\alpha+2\beta}&=-X_{3\alpha+\beta}.
    \end{aligned}
    \]
  \item[\SGAInlineRefTarget{sga3:XXIII:3.4.1:item:iii}(iii)] Put
    \[
    \begin{aligned}
      p_{\alpha+\beta}(x)
        &=\exp(xX_{\alpha+\beta})
          =\operatorname{int}(w_\beta)p_\alpha(x),\\
      p_{2\alpha+\beta}(x)
        &=\exp(xX_{2\alpha+\beta})
          =\operatorname{int}(w_\alpha w_\beta)p_\alpha(x),\\
      p_{3\alpha+\beta}(x)
        &=\exp(xX_{3\alpha+\beta})
          =\operatorname{int}(w_\alpha)p_\beta(-x),\\
      p_{3\alpha+2\beta}(x)
        &=\exp(xX_{3\alpha+2\beta})
          =\operatorname{int}(w_\beta w_\alpha)p_\beta(-x).
    \end{aligned}
    \]
    Then
    \[
    \begin{aligned}
      p_\beta(y)p_\alpha(x)
        ={}&p_\alpha(x)p_\beta(y)p_{\alpha+\beta}(xy)
          p_{2\alpha+\beta}(x^2y)p_{3\alpha+\beta}(x^3y)
          p_{3\alpha+2\beta}(x^3y^2),\\
      p_{\alpha+\beta}(y)p_\alpha(x)
        ={}&p_\alpha(x)p_{\alpha+\beta}(y)
          p_{2\alpha+\beta}(2xy)p_{3\alpha+\beta}(3x^2y)
          p_{3\alpha+2\beta}(3xy^2),\\
      p_{2\alpha+\beta}(y)p_\alpha(x)
        ={}&p_\alpha(x)p_{2\alpha+\beta}(y)
          p_{3\alpha+\beta}(3xy),\\
      p_{3\alpha+\beta}(y)p_\beta(x)
        ={}&p_\beta(x)p_{3\alpha+\beta}(y)
          p_{3\alpha+2\beta}(-xy),\\
      p_{2\alpha+\beta}(y)p_{\alpha+\beta}(x)
        ={}&p_{\alpha+\beta}(x)p_{2\alpha+\beta}(y)
          p_{3\alpha+2\beta}(3xy).
    \end{aligned}
    \]
\end{enumerate}}

\SGARefTarget{sga3:XXIII:3.4.2}
\SGAParagraph{3.4.2}
The proof occupies \SGARefLink{sga3:XXIII:3.4.2}{3.4.2}--\SGARefLink{sga3:XXIII:3.4.9}{3.4.9}. Since
\[
  (\beta^*,\alpha)=-1,\qquad(\alpha^*,\beta)=-3,
\]
we have \(\beta(t_\alpha)=\alpha(t_\beta)=-1\). Define the four
nonsimple root vectors as in \SGARefLink{sga3:XXIII:3.4.1:item:ii}{(ii)}. Then
\[
\begin{aligned}
  \operatorname{Ad}(w_\beta)X_{\alpha+\beta}
    &=\alpha(t_\beta)X_\alpha=-X_\alpha,\\
  \operatorname{Ad}(w_\alpha)X_{2\alpha+\beta}
    &=\alpha(t_\alpha)\beta(t_\alpha)X_{\alpha+\beta}
      =-X_{\alpha+\beta},\\
  \operatorname{Ad}(w_\alpha)X_{3\alpha+\beta}
    &=-\beta(t_\alpha)X_\beta=X_\beta,\\
  \operatorname{Ad}(w_\beta)X_{3\alpha+2\beta}
    &=\alpha(t_\beta)^3\beta(t_\beta)X_{3\alpha+\beta}
      =-X_{3\alpha+\beta}.
\end{aligned}
\]
Since \((3\alpha+2\beta)\pm\alpha\) and
\((2\alpha+\beta)\pm\beta\) are not roots,
\[
  \operatorname{Ad}(w_\alpha)X_{3\alpha+2\beta}=X_{3\alpha+2\beta},
  \qquad
  \operatorname{Ad}(w_\beta)X_{2\alpha+\beta}=X_{2\alpha+\beta}.
\]

\SGARefTarget{sga3:XXIII:3.4.3}
\SGAParagraph{3.4.3}
By \SGARefLink{sga3:xxii:cor:5.5.2}{Exposé XXII, 5.5.2}, there are sections
\[
  A,B,C,D,E,F,G,H,J\in\mathbf G_a(S)
\]
such that
\begin{align}
 p_\beta(y)p_\alpha(x)
  ={}&p_\alpha(x)p_\beta(y)p_{\alpha+\beta}(Axy)
     p_{2\alpha+\beta}(Bx^2y)p_{3\alpha+\beta}(Cx^3y)
     p_{3\alpha+2\beta}(Dx^3y^2), \tag{1}\label{sga3:XXIII:3.4:eq1}\\
 p_{\alpha+\beta}(y)p_\alpha(x)
  ={}&p_\alpha(x)p_{\alpha+\beta}(y)p_{2\alpha+\beta}(Exy)
     p_{3\alpha+\beta}(Fx^2y)p_{3\alpha+2\beta}(Gxy^2),
     \tag{2}\label{sga3:XXIII:3.4:eq2}\\
 p_{2\alpha+\beta}(y)p_\alpha(x)
  ={}&p_\alpha(x)p_{2\alpha+\beta}(y)p_{3\alpha+\beta}(Hxy),
     \tag{3}\label{sga3:XXIII:3.4:eq3}\\
 p_{3\alpha+\beta}(y)p_\beta(x)
  ={}&p_\beta(x)p_{3\alpha+\beta}(y)p_{3\alpha+2\beta}(Jxy).
     \tag{4}\label{sga3:XXIII:3.4:eq4}
\end{align}

\SGARefTarget{sga3:XXIII:3.4.4}
\SGAParagraph{3.4.4}
Conjugating (1), (3), and (4) by \(w_\beta\) gives
\begin{align}
 p_{-\beta}(-y)p_{\alpha+\beta}(x)
  ={}&p_{\alpha+\beta}(x)p_{-\beta}(-y)p_\alpha(-Axy)
     p_{2\alpha+\beta}(Bx^2y)\notag\\
    &{}\cdot p_{3\alpha+2\beta}(Cx^3y)
     p_{3\alpha+\beta}(-Dx^3y^2), \tag{5}\label{sga3:XXIII:3.4:eq5}\\
 p_{2\alpha+\beta}(y)p_{\alpha+\beta}(x)
  ={}&p_{\alpha+\beta}(x)p_{2\alpha+\beta}(y)
     p_{3\alpha+2\beta}(Hxy), \tag{6}\label{sga3:XXIII:3.4:eq6}\\
 p_{3\alpha+2\beta}(y)p_{-\beta}(-x)
  ={}&p_{-\beta}(-x)p_{3\alpha+2\beta}(y)
     p_{3\alpha+\beta}(-Jxy). \tag{7}\label{sga3:XXIII:3.4:eq7}
\end{align}
Conjugating (1) by \(w_\alpha\) gives
\begin{align}
 p_{3\alpha+\beta}(-y)p_{-\alpha}(-x)
  ={}&p_{-\alpha}(-x)p_{3\alpha+\beta}(-y)
     p_{2\alpha+\beta}(Axy)p_{\alpha+\beta}(-Bx^2y)\notag\\
    &{}\cdot p_\beta(Cx^3y)p_{3\alpha+\beta}(Dx^3y^2).
     \tag{8}\label{sga3:XXIII:3.4:eq8}
\end{align}

\SGARefTarget{sga3:XXIII:3.4.5}
\SGAParagraph{3.4.5}
The identity
\[
  w_\beta p_\alpha(x)w_\beta^{-1}=p_{\alpha+\beta}(x)
\]
may be written, since \(\alpha+2\beta\) is not a root,
\begin{equation}
  p_\beta(1)p_\alpha(x)p_\beta(-1)
   =p_{-\beta}(1)p_{\alpha+\beta}(x)p_{-\beta}(-1).
  \tag{9}\label{sga3:XXIII:3.4:eq9}
\end{equation}
Using (1) and then (4), its left side becomes
\begin{align}
 &p_\alpha(x)p_{\alpha+\beta}(Ax)p_{2\alpha+\beta}(Bx^2)
   p_{3\alpha+\beta}(Cx^3)
   p_{3\alpha+2\beta}((D-CJ)x^3). \tag{10}\label{sga3:XXIII:3.4:eq10}
\end{align}
Using (5) and then \SGARefLink{sga3:XXIII:3.4:eq7}{(7)}, its right side becomes
\[
 p_{\alpha+\beta}(x)p_\alpha(Ax)p_{2\alpha+\beta}(-Bx^2)
 p_{3\alpha+2\beta}(-Cx^3)p_{3\alpha+\beta}((CJ-D)x^3).
\]
Applying (2) transforms this into
\begin{align}
 &p_\alpha(Ax)p_{\alpha+\beta}(x)
 p_{2\alpha+\beta}((AE-B)x^2)
 p_{3\alpha+\beta}((A^2F+CJ-D)x^3)\notag\\
 &\hspace{35mm}\cdot
 p_{3\alpha+2\beta}((AG-C)x^3). \tag{12}\label{sga3:XXIII:3.4:eq12}
\end{align}
Comparing \SGARefLink{sga3:XXIII:3.4:eq10}{(10)} and \SGARefLink{sga3:XXIII:3.4:eq12}{(12)} gives
\[
  A=1,\qquad E=2B,\qquad C+D=F+CJ,\qquad F=G.
\]

\SGARefTarget{sga3:XXIII:3.4.6}
\SGAParagraph{3.4.6}
Write
\[
  w_\alpha p_\beta(y)w_\alpha^{-1}=p_{3\alpha+\beta}(-y).
\]
Since \(4\alpha+\beta\) is not a root,
\begin{equation}
 p_\alpha(1)p_\beta(y)p_\alpha(-1)
  =p_{-\alpha}(1)p_{3\alpha+\beta}(-y)p_{-\alpha}(-1).
 \tag{13}\label{sga3:XXIII:3.4:eq13}
\end{equation}
Using (1), the left side is
\[
 p_\beta(y)p_{\alpha+\beta}(-Ay)p_{2\alpha+\beta}(By)
 p_{3\alpha+\beta}(-Cy)p_{3\alpha+2\beta}(-Dy^2).
\]
Using \SGARefLink{sga3:XXIII:3.4:eq8}{(8)}, (6), and (4), the right side is
\[
\begin{aligned}
 p_\beta(Cy)p_{\alpha+\beta}(-By)p_{2\alpha+\beta}(Ay)
 p_{3\alpha+\beta}(-y)
 p_{3\alpha+2\beta}((D-CJ-ABH)y^2).
\end{aligned}
\]
Thus
\[
  C=1,\qquad A=B,\qquad D-CJ-ABH=-D.
\]
Together with the preceding identities, this yields
\[
 A=B=C=1,\quad E=2,\quad F=G,\quad
 D+1=F+J,\quad 2D=H+J.
\]

\SGARefTarget{sga3:XXIII:3.4.7}
\SGAParagraph{3.4.7}
From
\[
  w_\beta p_{3\alpha+\beta}(x)w_\beta^{-1}
    =p_{3\alpha+2\beta}(x)
\]
one obtains
\[
  p_{3\alpha+\beta}(x)p_{3\alpha+2\beta}(-Jx)
   =p_{3\alpha+2\beta}(x)p_{3\alpha+\beta}(-Jx),
\]
and therefore \(J=-1\).

\SGARefTarget{sga3:XXIII:3.4.8}
\SGAParagraph{3.4.8}
Finally,
\[
  w_\alpha p_{\alpha+\beta}(y)w_\alpha^{-1}=p_{2\alpha+\beta}(y),
\]
or equivalently
\[
 p_\alpha(1)p_{\alpha+\beta}(y)p_\alpha(-1)
  =p_{-\alpha}(1)p_\alpha(-1)p_{2\alpha+\beta}(y)
    p_\alpha(1)p_{-\alpha}(-1).
\]
Using (2) on the left and (3) on the right gives
\[
\begin{aligned}
 &p_{\alpha+\beta}(y)p_{2\alpha+\beta}(-Ey)
   p_{3\alpha+\beta}(Fy)p_{3\alpha+2\beta}(-Gy^2)\\
 &\qquad=p_{-\alpha}(1)p_{2\alpha+\beta}(y)
   p_{3\alpha+\beta}(Hy)p_{-\alpha}(-1).
\end{aligned}
\]
Commuting \(p_{-\alpha}(-1)\) first past
\(p_{3\alpha+\beta}(Hy)\), and then past \(p_{2\alpha+\beta}(y)\),
introduces no new \(p_{3\alpha+\beta}\)-factor. Comparing those
factors gives \(F=H\). Hence
\[
  2D=D+1,\qquad D=1,\qquad F=G=H=2D-J=3.
\]
This determines all the coefficients and proves \SGARefLink{sga3:XXIII:3.4.1:item:iii}{(iii)}.

\SGARefTarget{sga3:XXIII:3.4.9}
\SGAParagraph{3.4.9}
For \SGARefLink{sga3:XXIII:3.4.1:item:i}{(i)}, one obtains successively
\[
\begin{aligned}
 w_\alpha w_\beta w_\alpha
  &=w_{3\alpha+\beta}^{-1}t_\alpha,\\
 w_\beta(w_\alpha w_\beta)^2
  &=w_{3\alpha+2\beta}^{-1}t_\alpha,\\
 w_\alpha(w_\beta w_\alpha)^3
  &=w_{3\alpha+2\beta}^{-1},\\
 w_\beta(w_\alpha w_\beta)^4
  &=w_{3\alpha+\beta}t_\beta,\\
 w_\alpha(w_\beta w_\alpha)^5
  &=w_\beta^{-1}.
\end{aligned}
\]
Consequently
\[
  (w_\alpha w_\beta)^6=(w_\beta w_\alpha)^6=e.
\]

\SGARefTarget{sga3:XXIII:3.4.10}
\SGAStatement{Remark 3.4.10}{Condition \SGARefLink{sga3:XXIII:2.3:item:v}{(v)} of \SGARefLink{sga3:XXIII:2.3}{2.3} consists first of
\[
\text{(A)}\quad
\begin{cases}
 \operatorname{int}(h_\alpha h_\beta h_\alpha h_\beta h_\alpha)v_\beta
   =v_\beta,\\
 \operatorname{int}(h_\beta h_\alpha h_\beta h_\alpha h_\beta)v_\alpha
   =v_\alpha,
\end{cases}
\]
and, after setting
\[
\begin{aligned}
 v_{\alpha+\beta}&=\operatorname{int}(h_\beta)v_\alpha,&
 v_{2\alpha+\beta}&=\operatorname{int}(h_\alpha h_\beta)v_\alpha,\\
 v_{3\alpha+\beta}&=\operatorname{int}(h_\alpha)v_\beta^{-1},&
 v_{3\alpha+2\beta}&=\operatorname{int}(h_\beta h_\alpha)v_\beta^{-1},
\end{aligned}
\]
of the commutation relations
\[
\text{(B)}\quad
\begin{aligned}
 v_\beta v_\alpha
  &=v_\alpha v_\beta v_{\alpha+\beta}v_{2\alpha+\beta}
    v_{3\alpha+\beta}v_{3\alpha+2\beta},\\
 v_{\alpha+\beta}v_\alpha
  &=v_\alpha v_{\alpha+\beta}v_{2\alpha+\beta}^{\,2}
    v_{3\alpha+\beta}^{\,3}v_{3\alpha+2\beta}^{\,3},\\
 v_{2\alpha+\beta}v_\alpha
  &=v_\alpha v_{2\alpha+\beta}v_{3\alpha+\beta}^{\,3},\\
 v_{3\alpha+\beta}v_\alpha&=v_\alpha v_{3\alpha+\beta},&
 v_{3\alpha+2\beta}v_\alpha&=v_\alpha v_{3\alpha+2\beta},\\
 v_{\alpha+\beta}v_\beta&=v_\beta v_{\alpha+\beta},&
 v_{2\alpha+\beta}v_\beta&=v_\beta v_{2\alpha+\beta},\\
 v_{3\alpha+\beta}v_\beta
  &=v_\beta v_{3\alpha+\beta}v_{3\alpha+2\beta}^{-1},\\
 v_{3\alpha+2\beta}v_\beta&=v_\beta v_{3\alpha+2\beta},\\
 v_{2\alpha+\beta}v_{\alpha+\beta}
  &=v_{\alpha+\beta}v_{2\alpha+\beta}v_{3\alpha+2\beta}^{\,3},\\
 v_{3\alpha+\beta}v_{\alpha+\beta}
  &=v_{\alpha+\beta}v_{3\alpha+\beta},\\
 v_{3\alpha+2\beta}v_{\alpha+\beta}
  &=v_{\alpha+\beta}v_{3\alpha+2\beta},\\
 v_{3\alpha+\beta}v_{2\alpha+\beta}
  &=v_{2\alpha+\beta}v_{3\alpha+\beta},\\
 v_{3\alpha+2\beta}v_{2\alpha+\beta}
  &=v_{2\alpha+\beta}v_{3\alpha+2\beta},\\
 v_{3\alpha+2\beta}v_{3\alpha+\beta}
  &=v_{3\alpha+\beta}v_{3\alpha+2\beta}.
\end{aligned}
\]}
